% Frozen research candidate: protects a redesigned K11, not the original K12.
% The old printed 88-word PC2 is retained as the starting cycle C_old.
\paragraph{A derived three-row seed for $00001$.}
Define
\[
\begin{split}
D={}&\{12132\,13132,12132\,32132,12310\,32310,
13120\,23120,13123\,23123,13210\,23210,13230\,13232,\\
&21210\,21213,21230\,23230,21230\,31230,
21313\,21323,31212\,31232,32130\,32132\},\\
R={}&\{21210\,21213,21230\,31230\},\qquad K=D\setminus R.
\end{split}
\]
Thus $|D|=13$, $|R|=2$, $|K|=11$, and relative to the original
auxiliary set $K_0$,
\[
K=(K_0\setminus\{12313\,12323,13230\,23230\})
   \cup\{13120\,23120\}.
\]
For a listed base path $p=(v_0,\ldots,v_s)$, write $\partial p$
for its horizontal edges in both rows together with the two
vertical edges at $v_0,v_s$. Starting with the old printed two-row
cycle, successively toggle the edges of $\partial p$ for the seven
paths below. Every intermediate graph has degree two; its cycle
lengths, obtained by following the edges, are recorded on the right.
\[
\begin{array}{l|l}
p&\text{cycle lengths after the toggle}\\
\hline
(13120,23120,23123)&88\\
(12312,12313,12310)&39,49\\
(13213,13212,13232)&39,37,7,5\\
(12310,12320)&76,7,5\\
(13210,13212)&76,12\\
(13210,13230,23230)&72,16\\
(13132,13232)&88
\end{array}
\]
The final graph is therefore a two-row Hamilton cycle $C_2$.
It contains every edge of $D$ in both rows, and in the second row,
whose leading digit is $1$, it contains the matching
\[
M=\{12312\,12313,\;13212\,13213,\;13210\,13230\}.
\]
These three edges avoid $D$.

The following is a Hamilton path of the forty-state graph obtained
by deleting $12312,12313,13212,13213$; it contains every edge of $D$:
\begin{verbatim}
13210 23210 23213 21213 21210 21310 21313 21323 21320 31320
32320 32120 12120 13120 23120 23123 13123 12123 12323 12320
12310 32310 32312 31312 31310 31210 31212 31232 31230 21230
23230 23130 13130 12130 32130 32132 12132 13132 13232 13230
\end{verbatim}
Replace the three second-row edges of $M$ with paths through the
new third row, whose leading digit is $2$: use its edges
$12312\,12313$, $13212\,13213$ for the first two replacements and
the displayed forty-state path for the last one, adding a vertical
edge at each end. The replacement paths are internally disjoint
and cover the entire third row, so this subdivides three edges of
$C_2$ into a Hamilton cycle of all $132$ states. Since $M$ avoids
$D$ and the new paths retain $D$, all thirteen edges occur in
each of the three rows.

\paragraph{Interface changes.}
The six constant-source families retain their old witnesses:
their four canonical target edges
$31212\,31232$, $12132\,32132$, $21313\,21323$,
$13123\,23123$ all belong to $K$.
Of the 283 individual witnesses, exactly three require replacement:
\[
\begin{array}{c|c|c|c|c}
\text{source}&j&\text{new excluded colour}&\text{target}&\text{witness}\\
\hline
00001&1&3&00102&21313\,21323\\
00001&4&2&00001&13120\,23120\\
00001&4&3&00001&12310\,32310
\end{array}
\]
In the first row the prescribed target normalizer sends the witness
to $21231\,23231\in K_{00102}$. In the second row the target
colour transposition $(1\,2)$ exchanges $13120\,23120$ and
$13210\,23210$, both in $K$. In the third row the inverse target
map $(1,2,3)\mapsto(3,1,2)$ sends $13120\,23120$ to
$12310\,32310$. Thus the three new witnesses belong to both
specified auxiliary sets. The other 280 individual witnesses
retain their previous membership, as checked against the frozen
literal interface input.

% This candidate replaces the original K12 and its pool7 only for
% 00001. It leaves the constant-source families, the matching block
% chain in B00001, the Hamilton path in B00112, and all other types
% unchanged. No claim of a universal PC2-to-PC3 implication is made.
